\documentclass{article}

\begin{document}

\textbf{1.1 Nonlinear Cardiovascular System Examples}

\hfill


Before moving on to the main content of the book, a few cardiovascular examples
are provided so as to further motivate the study of Chaos and Fractals.

\hfill

\textbf{1.1.1 Blood Vessel Mechanics}

\hfill

Historically, researchers have used linear approximations to model physiological
systems. In most cases, this method has not presented a limitation to physiological
modeling since the first-order effects are represented adequately. However, the
Chaos response will not be represented since nonlinear properties are required to
observe a chaotic response. Blood vessel mechanics provides such an example. A
large difference between the results obtained via linear approximation versus a
complete solution of the nonlinear system occurs when not incorporating physiological
vessel mechanics. Consider the mechanics of a human brachial artery from
the perspective of its lumen area versus pressure relationship shown in Fig. 1.1. This
figure clearly demonstrates a nonlinear behavior, but to a large extent the vessel is
typically modeled as a constant compliance. Since the compliance is the slope of the
curve in Fig. 1.1, this consequently is not true. The linear compliance representation
of a blood vessel prevents any Chaos response from being observed by accurately
incorporating this effect in any analysis. Clearly, linear approximation of the “real”
arterial function would be limited to conditions where the normal physiology is
restrained to a narrow range of function. In this example, a blood vessel can easily
exhibit nonlinear function whenever the external stresses exceed that of the internal
pressure. This can occur when an intramuscular vessel, such as the heart, is subjected
to muscle contraction.

\hfill

\textbf{1.1.2 Cardiac Valve}


\hfill

Another example of important nonlinear function in the cardiovascular system is the
cardiac valve. The cardiac valve responds to the pressure across it by opening or
closing to allow or prevent blood flow. A positive pressure across the valve results
in aortic flow. The valve therefore functions as a highly nonlinear blood flow resistance.
Without the valve, the mechanics of cardiac pumping would render a very
inefficient process at moving blood in a single direction. The need for nonlinear
function is certain in the case of the valve. A hypothetical valve pressure–flow function
is shown in Fig. 1.2.

\hfill


Notice that blood flow is only permitted for positive values of valve pressure
drop. This creates a flow discontinuity about the zero pressure as well as a nonlinear
function for positive pressures. In the case of the valve, we find that its nonlinear function is an essential function of the cardiovascular system. Other dynamics that
arise from valve function are the heart sounds. The first and second heart sounds due
to the closure of the aortic and mitral valve are familiar to all cardiologists and are
a result of this nonlinear function. The full nonlinear impact of valvular function
remains to be studied regarding arterial blood pressure dynamics. It should be clear
at this point that nonlinear function is an essential element for Chaos to occur. In the
chapters that follow, it will be shown how nonlinear function leads to Chaos in a
system cardiovascular or otherwise. In physiology, Chaos will most often occur in
the variables that we observe, for example, in time signals. This is detected as a variability
in the signal. It might be assumed that the variation is due to a random process.
For example, the variation of a pure random process is shown in Fig. 1.3.

\hfill

\textbf{1.1.3 Chaotic Signals}

\hfill

Alternative to the random variation shown in Fig. 1.3, consider the variation of a nonlinear dynamic system that is generating chaos as shown in Fig. 1.4.

\hfill

The Chaos signal in Fig. 1.4 is a distinctly different kind of variability than that
of the random signal in Fig. 1.3. The random signal is truly random and unpredictable.
On the other hand, in this example, chaos was generated by a known system.
Hence, each point leads to the next in a deterministic way. Hence, we find a clear
value to our understanding of Chaos and nonlinear systems in that they are predictable
if the system that generated the Chaos is known at least in the short term. Our
conclusion from the random variable versus the chaos variable is quite different. In
the case of Chaos, the system may be defined. Referring again to Fig. 1.4, Chaos,
look for patterns in the time series. For example, you may find repeated values or
segments. This might be reasonable to expect since the data in Fig. 1.4 were generated
by a computer implementing an algorithm in a repeated manner. These repeated
time series patterns may be treated in a more organized way. That is the concept of
a Fractal. This leads to the second main topic of this book, which is the Fractal.

\hfill

\textbf{1.2 Fractals}

\hfill

As we have seen above, a time series may appear to have repeated values in it even though it is Chaos. This same kind of repetition is also found in geometry, particularly
in branching geometries.

\hfill

\textbf{1.2.1 Branching Structures}

\hfill


To begin the study of the Fractal concept, it will be necessary to examine repeated
structures such as branching. Biology favors the iterative process or repetition such
that a more complex system may be achieved. Branching is very common to much
of biology since it is the usual case that biology tends to duplicate one function repeatedly. This repetition or iteration can yield a very complex structure over time.
It is referred to here as a class of nonlinear function because it will be seen that most
iterative procedures yield complex nonlinear functions. Despite the resulting complex
process, it is generally possible to identify simplifying features of the resulting
process. For example, consider the process of blood vessel branching. Another
branching structure is that of the pulmonary airway system that can be created from
branching the trachea right down to the alveoli. Other organ systems may be constructed
in a similar fashion. For example, the branching structure of the arteries in
the cardiovascular system can be viewed as a similar branching structure. We recognize
these branching structures as Fractal geometry. Branching structure may also
be found in the electrophysiology of the heart. Referring to Fig. 1.5, the Purkinje
fiber system of the heart’s activation system is shown.

\hfill

Repetition also recalls biological oscillations, such as the pacemaker. Due to the
presence of nonlinearity in biological oscillators, they possess properties that are
distinct from the man-made oscillator. The characteristic properties of biological
oscillators will therefore be presented. Man-made physical oscillators are much
more rigid in the sense that they are designed to follow a math function such as the
sinewave. The pacemaker of the heart is of primary interest to the cardiologist. So,
the cardiac pacemaker will be reviewed as a nonlinear oscillator in this book. This
topic will be presented in the nonlinear cardiology chapter following this introduction
to some of the topics of this book.

\hfill

\textbf{1.3 Book Organization}

\hfill

This book has been organized into three main parts. Part I will introduce you to the
fundamentals of nonlinear dynamic systems and Chaos as well as methods to characterize
the nonlinear system response. Part II is devoted to an introduction to
Fractals. The Fractal will be introduced by way of using the concept to analyze
fractal objects and fractal time. In addition, the deterministic Fractal will be introduced
to provide a more mathematical foundation. Part III is devoted to applications
of what the reader has learned in Parts I and II primarily to the area of cardiovascular
dynamic research. Part III will show by example how to develop and analyze a
nonlinear dynamic model. Nonlinear methods will then be used to study the model
data. Lastly, this section will illustrate how the tools provided in this book can help
to eliminate assumptions and solve problems of cardiovascular physiology that have
otherwise gone unsolved using linear methods. The primary emphasis of the book
will be to supply the reader with first principles of Chaos and Fractals. This approach
will enable an understanding of the current research and the ability to solve research problems.

\hfill

\textbf{Problems}

\begin{enumerate}
\item List five human physiological systems that are nonlinear.
\item Sketch the nonlinear input/output function for each of the systems listed in
Problem 1.
\item Given that arterial compliance is dA/dP, sketch the arterial compliance curve for
the artery in Fig. 1.1.
\item Describe the frequency content of random noise.
\item What functional benefit does fractal geometry provide to the Purkinje fiber
network
of the heart?
\end{enumerate}

\hfill

\textbf{References}

\begin{enumerate}
\item Drzewiecki, G., \& Pilla, J. J. (1998). Noninvasive measurement of the human brachial artery
pressure-area relation in collapse and hypertension. Annals of Biomedical Engineering, 26(6),
965–974. https://doi.org/10.1114/1.130
\item Goldberger, A. L., \& West, B. J. (1987). Fractals in physiology and medicine. Yale Journal of
Biology and Medicine, 60(5), 421–435. Retrieved from https://www.ncbi.nlm.nih.gov/
pubmed/3424875
\end{enumerate}

\end{document}